\section{Kiến trúc hệ thống CI}

\subsection{Tổng quan kiến trúc}
Hệ thống CI được thiết kế với luồng hoạt động như sau:

\begin{enumerate}[noitemsep]
    \item Developer tạo branch mới, viết code và push lên GitHub
    \item Developer tạo Pull Request (PR) vào branch \texttt{main}
    \item GitHub Webhook tự động thông báo Jenkins
    \item Jenkins Multibranch Pipeline tự động phát hiện và chạy pipeline
    \item Pipeline thực thi 8 stages tuần tự
    \item Kết quả CI được báo về GitHub qua commit status
    \item PR chỉ được merge khi CI pass và có đủ 2 reviewer approve
\end{enumerate}

\subsection{Sơ đồ luồng CI}
Luồng CI pipeline bao gồm 8 stages chính:

\begin{table}[H]
\centering
\caption{8 Stages của CI Pipeline}
\begin{tabular}{|c|l|l|}
\hline
\textbf{Stage} & \textbf{Tên} & \textbf{Mô tả} \\
\hline
1 & Checkout & Clone code, fetch origin/main \\
\hline
2 & Detect Changed Modules & Xác định modules thay đổi (monorepo filter) \\
\hline
3 & Gitleaks -- Secret Scan & Quét secrets/API keys rò rỉ \\
\hline
4 & Test & Chạy unit tests + JaCoCo coverage report \\
\hline
5 & Coverage Gate & Kiểm tra coverage $\geq$ 70\% \\
\hline
6 & Build & Build các modules thay đổi \\
\hline
7 & SonarQube -- Analysis & Phân tích chất lượng code \\
\hline
8 & Snyk -- Dependency Scan & Quét lỗ hổng dependencies \\
\hline
\end{tabular}
\end{table}

\subsection{Hạ tầng và Kiến trúc Build phân tán (Master-Agent)}
\begin{itemize}[noitemsep]
    \item \textbf{Jenkins Server (Master node)}: Đóng vai trò làm máy chủ trung tâm (Orchestrator). Nhiệm vụ duy nhất là tiếp nhận tín hiệu từ GitHub Webhook, lập lịch, quản lý cấu hình Pipeline, lưu trữ lịch sử build và phân phối tác vụ xuống các Agent.
    \item \textbf{SonarQube Server}: Triển khai độc lập để phân tích chất lượng code của toàn bộ dự án.
    \item \textbf{Distributed Agents (Cộng tác viên build)}: Toàn bộ quá trình checkout code, build, unit test và phân tích bảo mật được giao hoàn toàn cho 3 Agent thực thi song song:
    \begin{itemize}[noitemsep]
        \item \textbf{Agent-Tri}: Phụ trách xử lý và thực thi các build job.
        \item \textbf{Agent-QVinh}: Phụ trách xử lý và thực thi các build job.
        \item \textbf{Agent-TVinh}: Phụ trách xử lý và thực thi các build job.
    \end{itemize}
    \item \textbf{Lợi ích}: Phân tán tải công việc giúp ngăn chặn tình trạng thắt nút cổ chai (bottleneck) tại Master node, đồng thời giúp tăng tốc thời gian hoàn thành pipeline lên gấp nhiều lần.
\end{itemize}
